% =============================================================
% eksaktalab-brand.tex — identitas visual EksaktaLab untuk worksheet
%
% Disisipkan oleh tools/build-worksheet.mjs tepat sebelum \begin{document}.
% Menimpa warna Matn* bawaan worksheet dan memasang header/footer EksaktaLab,
% sehingga file .tex sumber tidak perlu diubah.
%
% Warna mengikuti token di assets/css/materi.css (navy, teal, surface, border).
% \WsTopik diisi oleh build script (nama topik untuk header & footer).
% =============================================================

% --- Palet EksaktaLab ---
\definecolor{EkNavy}{HTML}{0A0F1E}
\definecolor{EkTeal}{HTML}{00BFA5}
\definecolor{EkTealDim}{HTML}{009E87}
\definecolor{EkTealSoft}{HTML}{E6F9F6}
\definecolor{EkSurface}{HTML}{F0F4FF}
\definecolor{EkBorder}{HTML}{D8E0F5}
\definecolor{EkText2}{HTML}{3D4A6B}
\definecolor{EkText3}{HTML}{7B8BB2}
\definecolor{EkAmber}{HTML}{F59E0B}
\definecolor{EkAmberSoft}{HTML}{FFF8E1}
\definecolor{EkBlue}{HTML}{3B82F6}
\definecolor{EkBlueSoft}{HTML}{EFF6FF}

% --- Petakan warna bawaan worksheet ke palet EksaktaLab ---
\colorlet{MatnPrimary}{EkNavy}        % strip judul, kotak indikator, kotak materi
\colorlet{MatnAccent}{EkTealDim}      % contoh terpandu, nomor soal
\colorlet{MatnLight}{EkSurface}       % latar kotak materi/rumus
\colorlet{MatnYellow}{EkTealSoft}     % latar contoh terpandu
\colorlet{MatnGreen}{EkAmberSoft}     % kotak "Ingat!" / catatan
\definecolor{MatnGray}{HTML}{F7F9FE}  % kotak jawaban
\colorlet{MatnBorder}{EkBorder}
\colorlet{MatnText}{EkNavy}
\colorlet{ShadeRegion}{EkBorder}
\colorlet{DrillBG}{EkBlueSoft}        % Mini Drill
\colorlet{DrillBorder}{EkBlue}
\colorlet{ThinkBG}{EkSurface}         % Think About It!
\colorlet{ThinkBorder}{EkNavy}
\colorlet{ScaffoldBG}{EkTealSoft}
\colorlet{ScaffoldBorder}{EkTealDim}

% --- Logo EksaktaLab: ikon (sama dengan assets/favicon.svg) + nama + tagline ---
% Satu tikzpicture agar posisinya presisi seperti navbar situs: nama rata atas ikon,
% tagline rata bawah ikon, keduanya rata kiri.
\newcommand{\EkLogoLockup}{%
  \begin{tikzpicture}[baseline=(icon.south)]
    \def\s{0.82cm}
    \node[inner sep=0, minimum size=\s, anchor=south west] (icon) at (0,0) {};
    \fill[EkTeal, rounded corners=0.2*\s] (icon.south west) rectangle (icon.north east);
    \begin{scope}[shift={(icon.north west)}, x=\s/448, y=-\s/448,
                  white, line width=0.085*\s, line cap=round, line join=round]
      % koordinat favicon (viewBox 512, kotak mulai di 32)
      \draw (128,168) -- (208,168);
      \draw (128,224) -- (228,224);
      \draw (128,280) -- (208,280);
      \draw (258,164) -- (318,224) -- (258,284);
    \end{scope}
    \node[inner sep=0, anchor=north west] at ([xshift=7pt, yshift=-3pt]icon.north east)
      {\headfont\bfseries\fontsize{15}{15}\selectfont\textcolor{EkNavy}{Eksakta}\textcolor{EkTeal}{Lab}};
    \node[inner sep=0, anchor=south west] at ([xshift=7.6pt, yshift=3pt]icon.south east)
      {\headfont\fontsize{6.2}{6.2}\selectfont\textcolor{EkText3}{\textls[140]{LAB BELAJAR MATEMATIKA}}};
  \end{tikzpicture}%
}

% --- Header & footer ---
\setlength{\headheight}{34pt}
\pagestyle{fancy}
\fancyhf{}
\fancyhead[L]{\EkLogoLockup}
\fancyhead[R]{\headfont\small\textcolor{EkText2}{Worksheet}\enspace
  \textcolor{EkTeal}{\textbullet}\enspace\textcolor{EkNavy}{\bfseries\WsTopik}}
\renewcommand{\headrulewidth}{0.8pt}
\renewcommand{\headrule}{{\color{EkTeal}\hrule width\headwidth height\headrulewidth}}

\fancyfoot[L]{\headfont\small\textcolor{EkText2}{\textcopyright\ \the\year\ \textcolor{EkNavy}{Eksakta}\textcolor{EkTealDim}{Lab}}}
\fancyfoot[C]{\headfont\textcolor{EkBorder}{\small\rule[0.5ex]{1.6cm}{0.4pt}\enspace
  \textcolor{EkTealDim}{\bfseries\thepage\ / \pageref{LastPage}}\enspace
  \rule[0.5ex]{1.6cm}{0.4pt}}}
\fancyfoot[R]{\headfont\small\textcolor{EkText2}{\WsTopik}}
\renewcommand{\footrulewidth}{0.4pt}
\renewcommand{\footrule}{{\color{EkBorder}\hrule width\headwidth height\footrulewidth}}
